% +++
% latex = "platex"
% +++
\RequirePackage{plautopatch}
\documentclass[dvipdfmx,a4paper]{jsarticle}
\usepackage{shokumu}

\title{Curriculum Vitae}
\author{Yume Ninomiya}
\date{Final Update Date: March 19, 2026}

% ページ番号を消す場合はコメントを外す
% \pagestyle{empty}

\begin{document}

\maketitle

\section*{\textbf{[ }Contact Information\textbf{ ]}}



\section*{\textbf{[}Education\textbf{ ]}}
% \begin{tabular}[h]{ll}
%     平成○年3月 & ○○大学大学院 ○○研究科 ○○専攻 博士課程修了 \\
%     & 【研究内容】○○におけるユーザ行動分析 \\
%     & 【実績】Journal of ○○掲載、○○学会発表 \\
% \end{tabular}

\section*{\textbf{[ }Employment History\textbf{ ]}}

% \begin{tabular}[h]{lp{12cm}}
%     平成○年○月 & 株式会社○○（本社○○市、従業員○人） \\
%      & 【事業概要】Webサービス運営 \\
%      & 【職務内容】自社サービスのマーケティング（アルバイト） \\
%      & 【実績】利用データより○○機能の利用が少ないことを発見した。そのため社内に○○機能への導線強化を提案した。○ヶ月後にMAU○○\%増を達成。 \\
% \end{tabular}

\section*{\textbf{[ }Presentations at Conferences and Seminars\textbf{ ]}}

% \paragraph*{データ分析・マーケティングスキル}

% 研究を通じて培ったデータ分析スキルがあります。またアルバイトを通じて経験した実務的なマーケティング能力もあります。
% そのため現場における生のデータを解析し、ビジネス価値につながる提案を行うことができます。
% なおPython、R、Googleアナリティクス、AWS、Illustrator、Photoshopなどを使うことができます。

% \paragraph*{挑戦することが好き}

% 何にでも新しいことに挑戦することが好きです。新しい技術を試したり、業界のトレンドにアンテナを張ることが好きです。体力も自信があります。

% \section{志望動機}

% 私は、データ分析や新しいことが好きです。そのためマーケティング職は、自分の天職だと思っております。このたび、今以上にもっとデータ分析に関わりたい、大きな仕事に挑戦したいと願い転職を希望しました。このような中で、企業戦略の刷新に対して提案する貴社から求人募集が出されたと知り、是非私も貴社の発展にお役に立ちたく思い応募しました。

\end{document}